% !TEX program = xelatex
% 数学建模竞赛论文模板（国赛风格）
% 编译：xelatex main.tex（需 TeX Live / MacTeX，含 ctex 宏包）
% 结构：摘要 → 一、问题重述 → 二、问题分析 → 三、模型假设 → 四、符号说明
%       → 五、模型建立与求解 → 六、模型检验与灵敏度 → 七、模型评价
%       → 参考文献 → 附录

\documentclass[12pt]{ctexart}

\usepackage[a4paper, top=2.5cm, bottom=2.5cm, left=2.5cm, right=2.5cm]{geometry}
\usepackage{amsmath, amssymb}
\usepackage{graphicx}
\usepackage{booktabs}      % 三线表
\usepackage{longtable}     % 跨页长表
\usepackage{multirow}
\usepackage{float}
\usepackage{caption}
\usepackage[hidelinks]{hyperref}

\linespread{1.4}
\setlength{\parindent}{2em}
\graphicspath{{../../work/figures/}}   % 按实际项目路径调整

\title{\textbf{题目：\LaTeX 论文模板}}
\author{参赛队号：XXXXXXXXXX}
\date{}

\begin{document}

\maketitle

% ==================== 摘要（单独一页） ====================
\begin{center}
{\Large\heiti 摘\quad 要}
\end{center}

本文针对\ldots\ldots 问题，建立了\ldots\ldots 模型。
针对问题一，采用\ldots\ldots 方法，得到\ldots\ldots（关键结果，含数字）。
针对问题二，\ldots\ldots。针对问题三，\ldots\ldots。

最后对模型进行了灵敏度分析与优缺点评价。

\noindent\textbf{关键词：}关键词一；关键词二；关键词三

\clearpage

% ==================== 一、问题重述 ====================
\section*{一、问题重述}

简要复述题目背景与要求，注意：\textbf{问题重述不是抄题}，要提炼要点。

% ==================== 二、问题分析 ====================
\section*{二、问题分析}

对每个小问逐一分析：需要什么数据、适合什么方法、如何组织求解。

% ==================== 三、模型假设 ====================
\section*{三、模型假设}

\begin{enumerate}
  \item 假设一（可辩护、可检验）。
  \item 假设二。
\end{enumerate}

% ==================== 四、符号说明 ====================
\section*{四、符号说明}

\begin{table}[H]
\centering
\caption{主要符号说明}
\begin{tabular}{clc}
\toprule
符号 & 含义 & 单位 \\
\midrule
$x_i$ & 第 $i$ 个决策变量 & —— \\
$f$   & 目标函数值 & —— \\
\bottomrule
\end{tabular}
\end{table}

% ==================== 五、模型建立与求解 ====================
\section*{五、模型建立与求解}

\subsection*{5.1 问题一模型的建立}

模型假设与符号说明基础上，建立如下模型：
\begin{equation}
  \min f = \sum_{i=1}^{n} c_i x_i
  \label{eq:obj}
\end{equation}
\begin{equation}
  \text{s.t.}\quad \sum_{i=1}^{n} a_{ij} x_i \le b_j, \quad j = 1, \dots, m
  \label{eq:cons}
\end{equation}

\subsection*{5.2 问题一模型的求解}

求解方法简述（求解器、算法、参数），给出结果并配图配表。

\begin{figure}[H]
\centering
\includegraphics[width=0.7\textwidth]{fig_1_result.png}
\caption{问题一结果示意图}
\end{figure}

\subsection*{5.3 问题一最终模型（迭代完成后的完整形式）}

% 硬性要求：给出迭代完成后的最终方程，全部参数数值代入
将参数代入方程 \eqref{eq:obj} 与 \eqref{eq:cons}，最终模型为：
\begin{equation}
  \min f = 0.32 x_1 + 0.58 x_2 + 0.41 x_3
\end{equation}

\subsection*{5.4 问题二……（同上结构）}

% ==================== 六、模型检验与灵敏度 ====================
\section*{六、模型检验与灵敏度分析}

误差指标（R² / RMSE / MAE 等）、交叉验证、残差分析、参数扰动 $\pm 5\% \sim \pm 20\%$ 的结果变化表。

\begin{table}[H]
\centering
\caption{灵敏度分析结果}
\begin{tabular}{lccc}
\toprule
参数 & $-20\%$ & 基准 & $+20\%$ \\
\midrule
$\alpha$ & 0.91 & 1.00 & 1.12 \\
\bottomrule
\end{tabular}
\end{table}

% ==================== 七、模型评价 ====================
\section*{七、模型评价（优缺点）}

客观列出模型优点与不足，以及改进方向。

% ==================== 参考文献 ====================
\begin{thebibliography}{9}
  \bibitem{ref1} 作者. 文章标题[J]. 期刊名, 年份, 卷(期): 页码.
  \bibitem{ref2} Author. Title[C]//Proceedings. Year.
\end{thebibliography}

% ==================== 附录 ====================
\section*{附录}

关键代码片段（完整代码见附件 \texttt{code\_appendix.pdf}）。

\end{document}
